% Generated draft: requires mathematical and source review.
\documentclass[11pt]{article}
\usepackage[margin=25mm]{geometry}
\usepackage{fontspec}
\setmainfont{Latin Modern Roman}
\usepackage{amsmath,amssymb,mathtools,unicode-math}
\setmathfont{Latin Modern Math}
\usepackage{enumitem,needspace,xurl,hyperref}
\setlength{\parindent}{0pt}
\setlength{\parskip}{5pt}
\setlength{\emergencystretch}{3em}
\newcommand{\N}{\mathbb{N}}
\newcommand{\Z}{\mathbb{Z}}
\newcommand{\Q}{\mathbb{Q}}
\newcommand{\R}{\mathbb{R}}
\newcommand{\C}{\mathbb{C}}
\newcommand{\F}{\mathbb{F}}
\newcommand{\A}{\mathbb{A}}
\newcommand{\PP}{\mathbb{P}}
\newcommand{\E}{\mathbb{E}}
\newcommand{\eps}{\varepsilon}
\newcommand{\dd}{\,\mathrm{d}}
\begin{document}
\section*{UnsolvedMath ID: 30002645}
\textbf{Source problem code:} catalogue entry
\subsection{Definitions and mathematical statement}
Let \(X\) be a \(k\)-cusped hyperbolic \(3\)-manifold, and let \(C(X)\) denote its set of cusps, so that \(|C(X)|=k\). Write \(\mathcal D(X)\) for the object called the deformation variety in the supplied record. The record does not specify its construction, its component, or the meaning of simplicity. Write \(\operatorname{SGI}_X(S,T)\) for the source's assertion that the cusps in \(S\) are strongly geometrically isolated from those in \(T\); this is a named, unresolved predicate because the record supplies no definition or deformation quantifiers.

The supplied conjectural assertion is that, for every such \(X\),
\[
\mathcal D(X)\text{ is not simple}
\quad\Longrightarrow\quad
\exists S\subseteq C(X)\;\operatorname{SGI}_X\bigl(S,C(X)\setminus S\bigr).
\]
The wording does not explicitly require \(\varnothing\ne S\ne C(X)\). Whether a nonempty proper subset is intended must be resolved before treating this display as a fully specified mathematical conjecture. Likewise, the allowed values of \(k\), completeness and finite-volume assumptions, orientability, the precise deformation variety, simplicity, and the scope and direction of strong geometric isolation remain unspecified. No definitions of these missing conditions are imposed here.

The supplied dated assessment describes a result of Jeon for holonomy varieties under rationally independent cusp shapes and related hypotheses. Those additional hypotheses and the relationship between the holonomy variety and \(\mathcal D(X)\) are not specified in the record. That assessment therefore does not provide a precise restricted theorem to append to this statement.
\subsection{Short English statement}
The source asserts that a cusped hyperbolic \(3\)-manifold whose deformation variety is not simple has a set of cusps strongly geometrically isolated from the remaining cusps. The defining predicates and the admissible cusp subsets require clarification. The supplied partially solved label is dated curation, not an independently verified conclusion.
\subsection{Sources}
\begin{itemize}
\item Local UnsolvedMath registry: supplied record with permanent ID \(30002645\), exact source problem code \texttt{catalogue entry}, title ``Nonsimple Deformation Varieties and Geometrically Isolated Cusps,'' in the Oberwolfach Reports Open Problems collection. No individual proposer is supplied; the proposed year is recorded as 2014.
\item Supplied report citation: \emph{Low-dimensional Topology and Number Theory} (2015), Oberwolfach Reports, pp.\ 2115--2176, DOI \texttt{10.4171/owr/2014/38}. Supplied literature-source label: ``original Oberwolfach report.'' \url{https://doi.org/10.4171/owr/2014/38}. The record describes Oberwolfach Reports as documenting mathematical workshops and their open-problem discussions.
\item Supplied literature-source label: ``Jeon, anomalous holonomy subvarieties.'' \url{https://arxiv.org/abs/2005.02481}. The supplied assessment says that Jeon proves the strong-isolation conclusion for holonomy varieties under rationally independent cusp shapes and related hypotheses, and that the unrestricted higher-cusped formulation in this record is not established. No exact theorem number or complete list of hypotheses is supplied.
\item Literature assessment date: \texttt{2026-08-21}. This is supplied dated curation derived from the OWR record, not independently verified here; the supplied background expressly requests independent verification of time-sensitive status and cited claims. The registry labels the problem \texttt{partially\_solved} and its statement \texttt{corrected\_verified}, with the explanation that mathematical content was preserved while notation, typography, or omitted definitions were repaired for publication. These are supplied registry assertions, not mathematical verification of this draft.
\end{itemize}
\end{document}
